\documentclass[12pt]{article}
\usepackage[utf8]{inputenc}
\usepackage[T1]{fontenc}
\usepackage[margin=1in]{geometry}
\usepackage{setspace}
\usepackage{xcolor}
\usepackage[colorlinks=true, urlcolor=blue!60!black, linkcolor=blue!60!black]{hyperref}
\setlength{\parindent}{0em}
\setlength{\parskip}{0.6em}
\renewcommand{\baselinestretch}{1.0}
\begin{document}
\begin{center}
{\LARGE \textit{United States v. Cruikshank}, 92 U.S. 542 (1876)\par}
\vspace{0.5em}
{\large Case Brief --- Two-Page Edition\par}
\end{center}

\section{Memory Jogger}
Reversing convictions from the Colfax Massacre, the Court held that the First and Second Amendments bind only Congress and that the Fourteenth Amendment reaches only state action (\href{https://supreme.justia.com/cases/federal/us/92/542/}{Justia}).

\section{Detailed Case Facts}
On Easter Sunday 1873, an armed white mob attacked Black freedmen guarding the Grant Parish courthouse in Colfax, Louisiana, killing scores after they surrendered. Federal prosecutors charged the attackers under Section~6 of the Enforcement Act of 1870 with conspiring to hinder victims' federal rights to assemble, bear arms, enjoy life and liberty, and vote.

\section{Procedural History}
Three defendants were convicted on sixteen conspiracy counts. Justice Bradley, on circuit, would have arrested judgment; Judge Woods disagreed, and the division was certified to the Supreme Court.

\section{Judicial Votes}
Arrest of judgment affirmed, 9--0 in result. Waite for eight Justices; Clifford concurred on pleading grounds only.

\section{Holding}
No count stated a federal offense. Congress may protect only rights of national citizenship. Assembly and arms rights restrict only Congress; the Fourteenth Amendment bars only state action; voting counts failed to allege racial motive; and the catch-all counts were fatally vague.

\section{Analysis of Opinions}
\textbf{Waite:} Under dual citizenship and \textit{Barron}, assembly and arms rights predate the Constitution and look to state protection. The Fourteenth Amendment guards only against States. Though the Court ``may suspect that race was the cause of the hostility,'' it was not averred. Indictments must ``descend to particulars.'' \textbf{Clifford:} the argued counts lacked particularity; he would not reach the constitutional questions.

\section{Examples: Future Applications}
\textbf{Same-Side:} A purely private assault with no federal nexus is for state law. \textbf{Opposite-Side:} Intimidating Black voters in a congressional election is federally punishable even under \textit{Cruikshank}. \textbf{Fence-Sitter:} A State's systematic refusal to prosecute racial terror may or may not be state action.

\section{Critique}
Originalists say \textit{Cruikshank} betrayed the Privileges or Immunities Clause's original meaning. Historians like James Gray Pope say it immunized white supremacist terror and opened the way to Jim Crow.

\section{Key Quotations}
\begin{quote}
``This is not a right granted by the Constitution. Neither is it in any manner dependent upon that instrument for its existence. The second amendment declares that it shall not be infringed, but this, as has been seen, means no more than that it shall not be infringed by Congress.''
\end{quote}
\noindent --- Chief Justice Waite, Opinion of the Court.

\end{document}
